%% kapa: exam=1 round=19 year=2008 period=1 form=A subject=civil_law subject_ko=민법 seq=24 number=24
%% source: Q-Net 감정평가사 기출문제 (cst00309) artlSeq=1001962 file=extracted/1008/1001962/2008년 제19회 감정평가사 1차시험 기출문제/1교시.hwp
%% conversion: kordoc 4.24.51 -> tools/gen1.py (auto); review=proofread-sonnet; flags=none
\begin{question}{24}
등기청구권에 관한 설명으로 틀린 것은? (다툼이 있는 경우에는 판례에 의함)
\begin{choices}
  \item 부동산의 매수인이 부동산을 인도받아 사용·수익하다가 그 부동산을 다른 사람에게 처분하고 그 점유를 승계하여 준 경우에도 이전등기청구권의 소멸시효는 진행되지 않는다.
  \item 매매계약이 합의해제된 경우 매도인의 매수인에 대한 소유권이전등기 말소청구권은 소멸시효에 걸린다.
  \item 소유권이전등기청구권을 매수인으로부터 양도받은 양수인은 매도인이 그 양도에 대하여 동의하지 않고 있다면 매도인에 대하여 채권양도를 원인으로 하여 소유권이전등기절차의 이행을 청구할 수 없다.
  \item 당사자간의 약정에 기초하여 환매권보류에 관한 등기청구권을 행사하는 경우, 그 청구권의 법적 성질은 채권적이다.
  \item 근저당권설정약정에 따른 근저당권설정등기청구권은 그 피담보채권이 되는 대여금채권과는 별개의 청구권으로서 시효기간 또한 독자적으로 진행된다.
\end{choices}
\end{question}
\answer{2}
